This section presents the central contribution: a mapping from the twelve technical
evaluation method families of Section~\ref{sec:taxonomy} to the requirements of the three
frameworks. We first state the coding scheme, then walk each framework, and close with a
cross-framework crosswalk. The full mapping is given in Table~\ref{tab:matrix} and
visualized in Fig.~\ref{fig:heatmap}.

\subsection{Technical-relevance coding}
For each (method family, requirement) pair we assign one of three levels of \emph{technical
relevance}, capturing how far the method's output could contribute to an evidentiary case for the
requirement, grounded in the official clause text. We code \emph{relevance}, not legal
sufficiency: a technical artifact rarely demonstrates conformity by itself, its adequacy
depending on intended purpose, applicable harmonised standards, acceptance criteria, and the
conformity-assessment route. \emph{Primary} (\direct): the output is itself the kind of
technical evidence the requirement chiefly asks for. \emph{Supporting} (\partialev): the
method contributes but is insufficient alone. \emph{Contextual} (\indirect): it informs the
requirement without serving as its evidence. A blank denotes non-applicability. Codes were
assigned in a first pass grounded in clause text, then challenged in a structured second
pass, with a second author independently coding a stratified sample (Section~\ref{sec:method}).
The coding is an interpretive engineering judgment (Section~\ref{sec:limits}); we present the
matrix as a \emph{proposed} crosswalk inviting external validation, not an established map of
what demonstrates conformity.

\paragraph{Worked example} To make the coding transparent, consider EU AI Act Article~15,
which requires high-risk systems to achieve ``appropriate levels of accuracy, robustness
and cybersecurity''~\cite{euaiact2024}. Performance and benchmarking (M1) has
\emph{primary} relevance to the accuracy dimension through held-out and out-of-distribution
test results; robustness and distribution-shift evaluation (M3) has \emph{primary} relevance
to robustness through adversarial and corruption tests; red-teaming (M5) has \emph{primary}
relevance to resilience to manipulation; uncertainty and calibration (M7) has only
\emph{supporting} relevance, since it characterizes the reliability of accuracy claims but
not adversarial security; and interpretability (M6) is not applicable to this
clause. Applying this logic clause by clause against the official text produces every cell
of Table~\ref{tab:matrix}.

% Auto-generated signature matrix (Table T4). Glyph macros defined in main.tex.
\begin{sidewaystable*}
\centering
\caption{Proposed methods-to-requirements crosswalk (the article\textquotesingle s central artifact): technical evaluation method families (rows) against regulatory requirements (columns). Each cell records the authors\textquotesingle{} interpretive coding of the method\textquotesingle s \emph{technical relevance} to the requirement: \direct~primary, \partialev~supporting, \indirect~contextual; a blank denotes not applicable. The codes are an engineering judgment grounded in the official clause text, not a legal determination of conformity; a second author independently coded a stratified 20\% sample (Section~\ref{sec:method}). NIST AI RMF functions/subcategories~\cite{nist2023airmf}, EU AI Act high-risk Articles~\cite{euaiact2024}, ISO/IEC 42001 Annex~A controls~\cite{iso42001}. ISO columns: A.5 assessing impacts, A.6.2.4 verification and validation, A.7 data for AI systems, A.8 information for interested parties, A.6.2.6 operation and monitoring, A.10 third-party relationships. Rotate the page 90$^\circ$ to read the column headers; Fig.~\ref{fig:heatmap} provides a visual overview of the same data.}
\label{tab:matrix}
\footnotesize
\setlength{\tabcolsep}{3pt}
\begin{tabular}{lcccccccccccccccccccccc}
\toprule
 & \multicolumn{9}{c}{\textbf{NIST AI RMF}} & \multicolumn{7}{c}{\textbf{EU AI Act}} & \multicolumn{6}{c}{\textbf{ISO/IEC 42001}}\\
\cmidrule(lr){2-10}\cmidrule(lr){11-17}\cmidrule(lr){18-23}
\textbf{Method family} & \rotatebox{90}{NIST GOVERN} & \rotatebox{90}{NIST MAP} & \rotatebox{90}{ME 2.3 Acc.} & \rotatebox{90}{ME 2.6 Safety} & \rotatebox{90}{ME 2.7 Sec.} & \rotatebox{90}{ME 2.9 Expl.} & \rotatebox{90}{ME 2.10 Priv.} & \rotatebox{90}{ME 2.11 Fair} & \rotatebox{90}{NIST MANAGE} & \rotatebox{90}{Art.9} & \rotatebox{90}{Art.10} & \rotatebox{90}{Art.11} & \rotatebox{90}{Art.12} & \rotatebox{90}{Art.13} & \rotatebox{90}{Art.14} & \rotatebox{90}{Art.15} & \rotatebox{90}{A.5 Impact} & \rotatebox{90}{A.6.2.4 V\&V} & \rotatebox{90}{A.7 Data} & \rotatebox{90}{A.8 Info} & \rotatebox{90}{A.6.2.6 Mon.} & \rotatebox{90}{A.10 3rd-pty}\\
\midrule
M1~Performance \& benchmarking & \indirect & \partialev & \direct & \indirect &  &  &  & \indirect & \partialev & \partialev & \indirect & \partialev &  & \partialev &  & \direct & \indirect & \direct & \indirect & \partialev & \indirect & \indirect\\
M2~Fairness \& bias evaluation & \indirect & \partialev & \indirect & \partialev &  &  &  & \direct & \partialev & \partialev & \partialev & \partialev &  & \partialev & \indirect & \partialev & \partialev & \partialev & \direct & \partialev & \partialev & \indirect\\
M3~Robustness \& distribution shift & \indirect & \partialev & \partialev & \partialev & \direct &  &  &  & \partialev & \partialev & \indirect & \partialev &  & \indirect &  & \direct & \indirect & \direct & \indirect & \indirect & \partialev & \\
M4~Safety evaluation & \indirect & \partialev & \indirect & \direct & \partialev &  & \indirect & \indirect & \partialev & \partialev &  & \partialev &  & \indirect & \partialev & \partialev & \partialev & \direct &  & \partialev & \partialev & \indirect\\
M5~Red-teaming \& adversarial testing & \indirect & \partialev &  & \partialev & \direct &  & \indirect & \indirect & \partialev & \partialev &  & \partialev &  &  & \indirect & \direct & \partialev & \direct &  & \indirect & \partialev & \indirect\\
M6~Interpretability \& explainability & \indirect & \indirect &  & \indirect &  & \direct &  & \partialev & \indirect & \indirect &  & \partialev &  & \direct & \partialev &  & \indirect & \partialev &  & \direct &  & \\
M7~Uncertainty \& calibration &  & \indirect & \partialev & \partialev &  & \indirect &  &  & \partialev & \partialev &  & \partialev &  & \partialev & \partialev & \partialev & \indirect & \partialev &  & \partialev & \partialev & \\
M8~Data quality, governance \& docs & \indirect & \direct & \indirect &  &  &  & \indirect & \partialev & \indirect & \indirect & \direct & \partialev &  & \partialev &  & \partialev & \partialev & \indirect & \direct & \partialev &  & \partialev\\
M9~Transparency \& model documentation & \indirect & \partialev & \indirect &  &  & \indirect & \indirect & \indirect & \indirect & \indirect & \indirect & \direct & \indirect & \direct & \partialev &  & \partialev & \indirect & \partialev & \direct & \indirect & \partialev\\
M10~Human oversight \& interaction & \indirect & \indirect &  & \partialev &  & \partialev &  & \indirect & \partialev & \indirect &  & \indirect &  & \partialev & \direct &  & \partialev & \partialev &  & \partialev & \partialev & \\
M11~Post-deployment monitoring \& assurance & \indirect & \indirect & \indirect & \partialev & \partialev &  & \indirect & \indirect & \partialev & \partialev &  & \indirect & \direct & \indirect & \indirect & \partialev & \partialev & \partialev & \indirect & \indirect & \direct & \partialev\\
M12~Privacy evaluation & \indirect & \partialev &  & \indirect & \partialev &  & \partialev & \indirect & \partialev & \partialev & \partialev & \partialev &  & \partialev &  & \partialev & \partialev & \partialev & \partialev & \indirect & \indirect & \indirect\\
\bottomrule
\end{tabular}
\end{sidewaystable*}


\subsection{Mapping to the NIST AI RMF}\label{sec:map-nist}
The NIST AI RMF concentrates technical measurement in the \textsc{Measure} function,
whose subcategories enumerate the trustworthiness characteristics to be
evaluated~\cite{nist2023airmf}. Throughout, we use NIST's official \textsc{Measure}~2.x
subcategory numbering (notably 2.3 performance, 2.5 validity and reliability, 2.6 safety,
2.7 security and resilience, 2.8 transparency and accountability, 2.9 explainability and
interpretability, 2.10 privacy, and 2.11 fairness and harmful bias managed). The matrix columns cover the
trustworthiness-characteristic subcategories most directly tied to technical evaluation;
the remaining subcategories (\textsc{Measure}~2.1, test
documentation; 2.2, human-subject evaluations; 2.4, in-production monitoring; 2.5,
validity and reliability; 2.8, transparency and accountability; 2.12, environmental
impact; and 2.13, measurement effectiveness) are addressed through the accuracy (2.3),
documentation (M8, M9), monitoring (M11), and governance rows rather than as separate
columns. The mapping is correspondingly densest in the \textsc{Measure} function.
Performance and benchmarking (M1) is primarily relevant to \textsc{Measure}~2.3
(accuracy/performance)~\cite{hendrycks2020mmlu,liang2023holistic}; robustness and
distribution-shift evaluation (M3) and red-teaming (M5) are primarily relevant to
\textsc{Measure}~2.7 (security and resilience) and supporting for \textsc{Measure}~2.6
(safety), for which dedicated safety evaluation (M4) is the primary
source~\cite{madry2018pgd,cohen2019certified,zou2023universal,mazeika2024harmbench};
interpretability (M6) is primarily relevant to \textsc{Measure}~2.9 (explainability and
interpretability)~\cite{lundberg2017shap,ribeiro2016lime}; and fairness evaluation (M2)
is primarily relevant to \textsc{Measure}~2.11 (fairness and harmful bias
managed)~\cite{parrish2022bbq,gallegos2024bias}; and privacy evaluation (M12)
has \emph{supporting} relevance to \textsc{Measure}~2.10 (privacy): its membership-inference,
extraction, and differential-privacy methods detect leakage and bound risk but do not by
themselves assure privacy adequacy~\cite{shokri2017membership,carlini2021extracting,abadi2016dp}.
Uncertainty and calibration (M7) has \emph{supporting} relevance to the accuracy and reliability
subcategories (\textsc{Measure}~2.3/2.5): it characterizes the reliability of performance
claims rather than producing the performance metric itself~\cite{guo2017calibration}. The
\textsc{Govern} and \textsc{Map} functions, being organizational and contextual, are served
only with \emph{supporting} or \emph{contextual} relevance: documentation families (M8, M9)
are the most relevant here~\cite{mitchell2019modelcards,gebru2018datasheets},
while \textsc{Manage}, covering ongoing risk treatment and monitoring, is chiefly served by
post-deployment monitoring and assurance (M11)~\cite{leest2025contextaware}. The pattern
is clear: NIST's measurement-centric design aligns tightly with the technical literature,
which is itself organized around measurable trustworthiness characteristics.

\subsection{Mapping to the EU AI Act (Articles 9--15)}\label{sec:map-eu}
The EU AI Act's high-risk obligations are more heterogeneous, mixing technical and
process requirements~\cite{euaiact2024}. Article~15 (accuracy, robustness, cybersecurity)
is the most technically saturated: M1 evidences accuracy, M3 robustness, and M5 the
resilience-to-manipulation and model-extraction concerns through adversarial testing and
prompt-injection
evaluation~\cite{dong2023promptrobust,greshake2023notwhat,zhang2024agent}. Article~10
(data and data governance) is primarily served by data-quality and documentation
methods (M8)~\cite{gebru2018datasheets,luccioni2023dataprovenance} and partially by
fairness evaluation of representativeness (M2), which evidences representativeness
outcomes rather than the data-governance management practices Article~10 requires. Article~11 (technical documentation, Annex~IV) and
Article~13 (transparency to deployers) are primarily served by the documentation and
transparency families (M9)~\cite{mitchell2019modelcards,pan2024modelreporting}. Article~14
(human oversight) is uniquely served by human-oversight evaluation
(M10)~\cite{guo2024reliance,chen2024reliancedrills}. Two obligations stand out for their
\emph{thin} technical coverage: Article~12 (record-keeping and logging) is an engineering
capability that few evaluation \emph{methods} assess for adequacy, and Article~9 (the
risk-management system) is a process that technical evaluations feed but do not
constitute. These observations anticipate the gap analysis.

\subsection{Mapping to ISO/IEC 42001}\label{sec:map-iso}
ISO/IEC~42001 is a management-system standard; its Annex~A controls are framed around
lifecycle processes rather than specific metrics~\cite{iso42001}. The technical
literature maps most directly onto the verification-and-validation control (A.6.2.4),
which nearly every method family can evidence, and onto the data controls (A.7), served
by M8. The AI-system impact-assessment controls (A.5) are substantially supported by
fairness and safety evaluation (M2, M4) but, as with EU Article~9, retain a process
character that technical methods only partly satisfy, so those cells are coded
supporting rather than primary. The information-to-interested-parties control
(A.8) is served by transparency documentation (M9), and operation-and-monitoring
(A.6.2.6) by M11. Because ISO/IEC~42001 is explicitly designed to interoperate with risk
frameworks~\cite{iso23894} and is referenced by the EU AI Act's quality-management
obligation, evidence assembled for one framework substantially transfers to it.

\subsection{Cross-framework crosswalk}\label{sec:crosswalk}
A striking regularity emerges across the three mappings: the same technical evaluation
artifact frequently evidences obligations in all three frameworks simultaneously
(Fig.~\ref{fig:crosswalk}). An adversarial-robustness report evidences NIST
\textsc{Measure}~2.7, EU Article~15, and ISO/IEC~42001 A.6.2.4 at once; a model card
evidences NIST documentation expectations, EU Articles~11 and~13, and ISO/IEC~42001 A.8.
This convergence is the practical engine behind the practitioner playbook of
Section~\ref{sec:discussion}: a single, well-designed evaluation dossier can serve
multiple jurisdictions, provided it is explicitly cross-referenced to each framework's
clauses. We stress that this convergence is an \emph{interpretive} claim of evidentiary
compatibility, namely that the same artifact can satisfy the evidentiary need of each clause, and
not an assertion that the obligations are legally identical; their scope, thresholds, and
legal force differ, and the alignments in Fig.~\ref{fig:crosswalk} should be read as
reasoned correspondences rather than equivalences. The mapping matrix is, in effect, the
index that makes such reuse auditable.
